%% Philipp G Freimann Juli 2019 für die BBW
%% Phi BBW-Vorlage für Mathematische Dokumente (LaTeX)
%% 2019 - 07 - 11

%%  In den Dokumenten können die folgenden Attribute überschrieben werden:

%%%%%%%%%%%%%%%%%%%%%%% P A C K A G E S %%%%%%%%%%%%%%%%%%%%%%%%%%%%%
\documentclass[twoside,14pt,a5paper]{extarticle}
\usepackage[paper=a4paper,margin=17mm]{geometry}%

\input{inputs/LaFLInput_Standalone}

\usepackage{FL_LayoutStandaloneSty}

%%%%%%%%%%%%%%%  H E A D E R   &   F O O T E R %%%%%%%%%%%%%%%%%%%%

%% Oben (Header) linke Seite
%\fancyhf[HLE]{\makebox{\includegraphics[width=37mm]{logos/bbwBreit.pdf}}} 
\fancyhf[HCE]{\parttitle}
%% Oben (Header) rechte Seite
\fancyhf[HRO]{\leftmark}
%% Unten (Footer)  FRE = right even, FLE= left even, FRO = right odd,
%% FCO = center odd
\fancyhf[FRE]{\doctitel{}:\ \fachthema}
\fancyhf[FLE,FRO]{\thepage{}/\pageref{LastPage}}
\fancyhf[FCO]{KST}
